\documentclass[11pt]{article}
\usepackage[margin=0.86in]{geometry}
\usepackage[T1]{fontenc}
\usepackage{lmodern}
\usepackage{microtype}
\usepackage{amsmath,amssymb,amsthm,mathtools}
\usepackage{xcolor}
\usepackage{graphicx}
\usepackage[colorlinks=true,linkcolor=blue!55!black,urlcolor=blue!55!black,citecolor=blue!55!black]{hyperref}

\newtheorem{theorem}{Theorem}
\newcommand{\M}{\mathcal M}
\newcommand{\T}{\mathcal T}
\newcommand{\one}{\mathbf 1}
\newcommand{\N}{\mathbb N}
\newcommand{\dens}{\operatorname{dens}}

\title{A Bounded-Generator Counterexample to\\
Semigroup Little-H\"older Biduality}
\author{Automated mathematical discovery packet for arXiv:1904.10332\\
\small model: GPT5.6; lane: agent\_lane\_08}
\date{19 July 2026}

\begin{document}
\maketitle

\begin{abstract}
Gonz\'alez-P\'erez asked whether the little semigroup H\"older class
$\lambda_\alpha(\T)$ always has bidual isomorphic to the big class
$\Lambda_\alpha(\T)$.  We give a negative answer for every
$0<\alpha<1$, already for an ergodic symmetric Markov semigroup on the
finite commutative von Neumann algebra $\ell_\infty(\N)$.  The generator is
the bounded projection $A=I-E$, where $E$ is expectation onto the constants.
Consequently all elements are little H\"older and both classes equal
$\ell_\infty$ with equivalent norms.  Nevertheless
$\ell_\infty^{**}$ has strictly larger density character than
$\ell_\infty$, so the two spaces are not even abstractly isomorphic.
\end{abstract}

\section{Source question}

For a Markov semigroup $\T=(T_t)_{t\geq0}$ with generator $A$ and
subordinated Poisson semigroup $P_s=e^{-sA^{1/2}}$, the source defines
\begin{align*}
 \Lambda_\alpha(\T)
 &=\left\{f\in\M:
   \sup_{s>0}\left\|s^{1-\alpha}\frac{dP_s}{ds}(f)\right\|_\infty
   <\infty\right\},\\
 \|f\|_{\Lambda_\alpha}
 &=\max\left\{\|f\|_\infty,
   \sup_{s>0}\left\|s^{1-\alpha}\frac{dP_s}{ds}(f)\right\|_\infty
   \right\}.
\end{align*}
Open Problem P.3 on page 34 of the arXiv PDF (Section 8.1, page 3282 of the
published version) defines
\[
 \lambda_\alpha(\T)=\left\{f\in\Lambda_\alpha(\T):
 \lim_{s\to0^+}\left\|s^{1-\alpha}\frac{dP_s}{ds}(f)\right\|_\infty=0
 \right\}
\]
and asks whether
$\lambda_\alpha(\T)^{**}\cong\Lambda_\alpha(\T)$ \cite{GonzalezPerez}.

\begin{center}
\includegraphics[width=0.98\textwidth]{figures/open_problem_crop.png}
\end{center}

\section{Proof intuition}

The proposed little/big distinction only detects singular short-time
behaviour of the generator.  If $A^{1/2}$ is bounded, then
$s^{1-\alpha}A^{1/2}P_s f\to0$ uniformly for every $f$, so the little and big
classes collapse to the ambient von Neumann algebra.  Choosing the ambient
algebra to be $\ell_\infty$ turns this collapse into an obstruction: the
desired biduality would assert that $\ell_\infty^{**}$ is isomorphic to
$\ell_\infty$.  A quotient-density argument rules out even an arbitrary
Banach-space isomorphism.  The expectation semigroup below is finite,
commutative, symmetric, ergodic, and nonconstant, so the obstruction is not
caused by the identity semigroup.

\begin{theorem}[Counterexample to P.3]\label{thm:main}
For every $0<\alpha<1$ there are a finite commutative von Neumann algebra
$(\M,\tau)$ and an ergodic symmetric Markov semigroup $\T$ on $\M$ such that
\[
 \lambda_\alpha(\T)=\Lambda_\alpha(\T)=\M
\]
with equivalent norms, but
\[
 \lambda_\alpha(\T)^{**}\not\cong\Lambda_\alpha(\T).
\]
Thus Open Problem P.3 has a negative answer in its stated generality.
\end{theorem}

\section{Proof}

Let $\M=\ell_\infty(\N)$ and equip it with the faithful normal finite trace
\[
 \tau(f)=\sum_{n=1}^{\infty}2^{-n}f_n,
 \qquad f\in\M_+.
\]
Since $\M$ is commutative, this normal state is a trace.  Let
\[
 E(f)=\tau(f)\one
\]
be the trace-preserving conditional expectation onto the constants and put
\begin{equation}\label{eq:semigroup}
 T_t=E+e^{-t}(I-E)=e^{-t}I+(1-e^{-t})E,
 \qquad t\geq0.
\end{equation}
The map $E$ is normal, unital, completely positive, trace preserving, and an
orthogonal projection on $L^2(\M,\tau)$.  Hence every $T_t$ is Markovian and
self-adjoint on $L^2$.  Since $E^2=E$, formula \eqref{eq:semigroup} also gives
$T_tT_u=T_{t+u}$ and pointwise weak-* continuity.  Its fixed-point algebra is
$\mathbb C\one$, so the semigroup is ergodic.  Its positive self-adjoint
generator is
\[
 A=I-E.
\]
This is an orthogonal projection, whence $A^{1/2}=A$.  Therefore the Poisson
semigroup equals $T_s$ and
\begin{equation}\label{eq:derivative}
 P_s=E+e^{-s}(I-E),
 \qquad
 \frac{dP_s}{ds}=-e^{-s}(I-E).
\end{equation}

Fix $0<\alpha<1$.  It follows from \eqref{eq:derivative} that for every
$f\in\M$,
\begin{align}
 \sup_{s>0}\left\|s^{1-\alpha}\frac{dP_s}{ds}(f)\right\|_\infty
 &=c_\alpha\,\|(I-E)f\|_\infty,\label{eq:seminorm}\\
 c_\alpha&=\sup_{s>0}s^{1-\alpha}e^{-s}
 =(1-\alpha)^{1-\alpha}e^{-(1-\alpha)}.
\end{align}
In particular, $\Lambda_\alpha(\T)=\M$.  Moreover,
\[
 \lim_{s\to0^+}s^{1-\alpha}e^{-s}\|(I-E)f\|_\infty=0,
\]
so every element is little H\"older and
$\lambda_\alpha(\T)=\Lambda_\alpha(\T)=\M$.  Their common inherited norm is
\[
 \|f\|_{\Lambda_\alpha}
 =\max\{\|f\|_\infty,c_\alpha\|(I-E)f\|_\infty\},
\]
which is equivalent to $\|f\|_\infty$ because $E$ is contractive and hence
$\|(I-E)f\|_\infty\leq2\|f\|_\infty$.

It remains to show that $\M^{**}$ is not isomorphic to $\M$.  Write
$\mathfrak c=2^{\aleph_0}$.  The set $\{0,1\}^{\N}\subset\ell_\infty$ is
1-separated and has cardinal $\mathfrak c$, while
$|\ell_\infty|=\mathfrak c$.  Thus
\begin{equation}\label{eq:density-M}
 \dens(\M)=\mathfrak c.
\end{equation}

We next exhibit a much larger quotient of $\M^{**}$.  Identify
$\ell_\infty=C(\beta\N)$.  The compactification $\beta\N$ contains a set
$I$ of $\mathfrak c$ distinct points.  For completeness, choose a family of
$\mathfrak c$ pairwise almost-disjoint infinite subsets of $\N$ (encode the
finite initial segments of each infinite binary sequence) and extend each to
a free ultrafilter; the resulting ultrafilters are distinct.  The point
masses $(\delta_p)_{p\in I}$ span an isometric copy
\[
 Y=\ell_1(I)\subset C(\beta\N)^*=\M^*.
\]
The restriction map
\[
 R:\M^{**}=(\M^*)^*\longrightarrow Y^*=\ell_\infty(I)
\]
is surjective by Hahn--Banach.  Density cannot increase under a continuous
surjection.  On the other hand, $\{0,1\}^{I}\subset\ell_\infty(I)$ is
1-separated and has cardinal $2^{\mathfrak c}$.  Consequently
\begin{equation}\label{eq:density-bidual}
 \dens(\M^{**})\geq\dens(\ell_\infty(I))\geq2^{\mathfrak c}
 >\mathfrak c=\dens(\M).
\end{equation}
Since density character is an isomorphic invariant, \eqref{eq:density-M} and
\eqref{eq:density-bidual} imply $\M^{**}\not\cong\M$.  Together with the
identification of the two H\"older classes above, this proves the theorem.

\section{Verification}

From this packet directory, run
\begin{verbatim}
conda run --no-capture-output -n sandbox python \
  code/verify_expectation_semigroup.py
\end{verbatim}
The script uses exact rational arithmetic on finite weighted analogues of
$E$ for $N=2,3,5,8,16$.  It verifies unitality, trace preservation,
$L^2(\tau)$ self-adjointness, $E^2=E$, $A^2=A$, and the exact semigroup
product formula for several rational parameters.  It also reports the
closed-form maximizing point $s=1-\alpha$ in \eqref{eq:seminorm}.  These checks
are sanity tests for the explicit algebra; the infinite-dimensional density
argument is the formal proof and is not replaced by computation.
The recorded run ended with \texttt{all checks passed}; the complete output is
stored in \texttt{verification\_report.txt}.

\section{Novelty bounds, limitations, and review recommendation}

A bounded search on 19 July 2026 covered the run's four lightweight indexes,
the exact arXiv id and title, the exact P.3 formula and phrases ``little
H\"older'', ``semigroup H\"older bidual'', and ``bounded generator'', the
published 2022 article, and later semigroup Campanato literature.  The closest
later item found was Hong--Jing's 2024 paper \cite{HongJing}, which answers a
Campanato/Lipschitz equivalence problem but does not discuss P.3 or biduality.
No source containing this counterexample or an answer to P.3 was found.
Novelty confidence is moderate--high, subject to a full citation review; the
construction is elementary enough that an uncatalogued observation remains
possible.

The theorem refutes only the unrestricted formulation of P.3.  A positive
version may still hold under hypotheses forcing genuine short-time smoothing,
or otherwise ensuring that $\lambda_\alpha(\T)$ is a proper subspace of
$\Lambda_\alpha(\T)$.  It does not address P.1, P.2, P.4, or P.5.

Human review is recommended, with focus on the interpretation of the source's
Banach-space isomorphism, the density-quotient step, and whether an intended
unstated hypothesis was meant to exclude bounded generators.  The density
argument deliberately proves failure of arbitrary isomorphism, not merely
failure of the canonical bidual map.

\begin{thebibliography}{9}
\bibitem{GonzalezPerez}
A. Gonz\'alez-P\'erez,
\emph{H\"older classes via semigroups and Riesz transforms},
Math. Z. 301 (2022), 3245--3284;
\href{https://arxiv.org/abs/1904.10332}{arXiv:1904.10332}.

\bibitem{HongJing}
G. Hong and Y. Jing,
\emph{Campanato spaces via quantum Markov semigroups on finite von Neumann
algebras},
Int. Math. Res. Not. IMRN 2024, no. 21, 13611--13641;
\href{https://arxiv.org/abs/2409.14414}{arXiv:2409.14414}.

\bibitem{HindmanStrauss}
N. Hindman and D. Strauss,
\emph{Algebra in the Stone--\v{C}ech Compactification}, 2nd ed.,
de Gruyter, 2012.
\end{thebibliography}

\end{document}
